\section{Massa va energiya orasidagi bog‘lanish}


\begin{enumerate}

\item Tinchlikdagi massasi 20 kg bo‘lgan jismning tinchlikdagi energiyasi necha Joulga teng?
\begin{enumerate}[label=\Alph*)]
\item $2,4\cdot 10^{20}$.
\item $1,2\cdot 10^{18}$.
\item $1,8\cdot 10^{18}$.
\item $0,9\cdot 10^{20}$.
\end{enumerate}

\item Tinchlikdagi massasi 14 kg bo‘lgan jismning tinchlikdagi energiyasi necha Joulga teng?
\begin{enumerate}[label=\Alph*)]
\item $1,3\cdot 10^{15}$.
\item $2,33\cdot 10^{14}$.
\item $1,82\cdot 10^{19}$.
\item $1,26\cdot 10^{18}$.
\end{enumerate}


\item Qanday tezlikda zarrachaning kinetik energiyasi uning tinchlikdagi energiyasiga teng bo‘ladi (m/s)?
\begin{enumerate}[label=\Alph*)]
\item $1,3\cdot 10^8$.
\item $1,3\cdot 10^7$.
\item $2,6\cdot 10^7$.
\item $2,6\cdot 10^8$.
\end{enumerate}

\item Tinchlikdagi massasi $m_0$ bo‘lgan jism $0,6c$ tezlik bilan harakatlansa to‘la energiyasi nimaga teng bo‘ladi?
\begin{enumerate}[label=\Alph*)]
\item $2,5m_0c^2$.
\item $3m_0c^2$.
\item $1,8m_0c^2$.
\item $1,25m_0c^2$.
\end{enumerate}

\item Zarraning tezligi $0,8c$ ga teng. Uning kinetik energiyasining to‘liq energiyasiga nisbati topilsin. $c$ -- yorug‘likning vakuumdagi tezligi.
\begin{enumerate}[label=\Alph*)]
\item 2/5.
\item 1/5.
\item 2/3.
\item 1/4.
\end{enumerate}

\item Qanday tezlikdagi zarrachaning kinetik energiyasi uning tinchlikdagi energiyasidan 2 marta katta bo‘ladi? $c$ -- yorug‘likning vakuumdagi tezligi.
\begin{enumerate}[label=\Alph*)]
\item $\sqrt{3}c/3$.
\item $2\sqrt{2}c/3$.
\item $c/3$.
\item $\sqrt{2}c/3$.
\end{enumerate}

\item Protonning kinetik energiyasi $0,5m_0c^2$ bo‘lsa, uning to‘la energiyasi qanday? $m_0$ -- protonning tinchlikdagi massasi.
\begin{enumerate}[label=\Alph*)]
\item $1,5m_0c^2$.
\item $0,5m_0c^2$.
\item $m_0c^2$.
\item $2m_0c^2$.
\end{enumerate}


\item Suv muzlaganda massasi o‘zgaradimi?
\begin{enumerate}[label=\Alph*)]
\item kamayadi.
\item ortadi.
\item o‘zgarmaydi.
\item suvning massasiga bog‘liq.
\end{enumerate}

\item 1,8 kg suv $45^\circ\text{C}$ ga isitilganda, uning massasi qanchaga ortadi (kg)? $c=4200\text{ J}/(\text{kg}\cdot\text{K})$.
\begin{enumerate}[label=\Alph*)]
\item $3,78\cdot 10^{-12}$.
\item $2,8\cdot 10^{-13}$.
\item $8,1\cdot 10^{-13}$.
\item $3,78\cdot 10^{-13}$.
\end{enumerate}

\item 3 kg suv $30^\circ\text{C}$ isitilganda, massasi qanday o‘zgaradi? Suvning solishtirma issiqlik sig‘imi $4200\text{ J}/(\text{kg}\cdot\text{K})$.
\begin{enumerate}[label=\Alph*)]
\item o‘zgarmaydi.
\item $42\cdot 10^{-13}$ kg ga ortadi.
\item $42\cdot 10^{-13}$ kg ga kamayadi.
\item $84\cdot 10^{-13}$ kg ga ortadi.
\end{enumerate}

\item Massasi 20 kg bo‘lgan azot o‘zgarmas bosimda $0^\circ\text{C}$ dan $200^\circ\text{C}$ gacha qizdirildi. Bunda azotning massasi qanchaga oshgan? Azotning o‘zgarmas bosimdagi solishtirma issiqlik sig‘imi $1,05\text{ kJ}/(\text{kg}\cdot\text{K})$.
\begin{enumerate}[label=\Alph*)]
\item $4,7\cdot 10^{-5}$ g.
\item $4,7\cdot 10^{-11}$ g.
\item 0.
\item $4,7\cdot 10^{-8}$ g.
\end{enumerate}

\end{enumerate}


\begin{table}[!ht]
\centering
\begin{tabular}{|c|c|c|c|c|c|c|c|c|c|c|}
\hline
1c & 2d & 3d & 4d & 5a & 6b & 7a & 8c & 9a & 10c & 11d \\ \hline
\end{tabular}
\end{table}